% GDZ PPN236018647,native543middle–544before61;HKUST544–545.
\NumberedItem{59}
\indent Formularum integralium formae:
\[
\int e^{-t^\alpha}dt\quad\text{et}\quad\int\frac{du}{{}^{\beta}\!\surd(1+u^\gamma)}
\]
\noindent inter se comparationem institui.
\par\noindent\hfill[1797] M[a]rt. 2.

\medskip\begingroup\small\makeatletter\let\veqno\@@leqno\makeatother
\indent Wahrscheinlich sind hier die zwischen den Grenzen \(0\) und \(+\infty\) genommenen bestimmten Integrale gemeint. Die Substitutionen
\[
t^\alpha=x\quad\text{beziehungsweise}\quad u^\gamma=\frac{x}{1-x}
\]
\noindent ergeben nämlich
\begin{equation*}\tag*{(1)}
\int_0^{+\infty}e^{-t^\alpha}dt=\frac1\alpha\int_0^{+\infty}e^{-x}x^{\frac1\alpha-1}dx,
\end{equation*}
\begin{equation*}\tag*{(2)}
\int_0^{+\infty}\frac{du}{\sqrt[\beta]{1+u^\gamma}}=\frac1\gamma\int_0^1x^{\frac1\gamma-1}(1-x)^{\frac1\beta-\frac1\gamma-1}dx,
\end{equation*}
\par\endgroup

\clearpage\SourceLeaf{544}\ArithmeticHead{515}{R}{1797 FEBR. 16—1797 MART.}
\begingroup\small
\noindent so daß die Beziehung zwischen den Integralen (1) und (2) auf die bekannte Gleichung zwischen den \textsc{Euler}schen Integralen hinauskommt, die \textsc{Euler} in der Abhandlung \textit{Evolutio formulae integralis \(\int x^{f-1}dx\,(lx)^{\frac mn}\) integratione a valore \(x=0\) ad \(x=1\) extensa}\footnote{Novi Commentarii Acad. Petrop. \textbf{16} (1771), 1772, S. 91, L. \textsc{Euleri} Opera omnia, ser. I, vol. 17, S. 316, siehe insbesondere § 25, S. 330 und das \textit{Supplementum}, S. 354 der Opera.} angegeben hat, und die in den jetzt üblichen Bezeichnungen so lautet:
\[
\mathrm{B}(p,q)=\frac{\Gamma(p)\Gamma(q)}{\Gamma(p+q)};
\]
\noindent vergl. auch den art. 28. der \textsc{Gauss}schen Abhandlung \textit{circa seriem} vom Jahre 1812, Werke III, S. 251.
\par\noindent\hfill\textsc{Schlesinger.}
\par\endgroup

\NumberedItem{60}
\indent Cur ad aequationem perveniatur gradus \(nn^{\mathrm{ti}}\) dividendo curvam lemniscatam in \(n\) partes.
\par\noindent\hfill[1797] M[a]rt. 19.

\medskip\begingroup\small
\indent Während der Grad der Gleichung, von der die Teilung des Kreises in \(n\) gleiche Teile abhängt, gleich der Anzahl der Teile ist, so ergab sich für die mit Hilfe des \textsc{Euler}schen Additionstheorems gebildete Gleichung, von der die Teilung der Lemniskate in \(n\) Teile abhängt, der Grad gleich \(n^2\). Der Grund hierfür, d. h. die Bedeutung nicht nur der \(n\) realen, sondern auch der \(n^2-n\) komplexen Wurzeln\begingroup\renewcommand{\thefootnote}{**)}\footnote{Vergl. oben S. 162 das »Determ. rad. imag.« für die Fünfteilung. Im Text der Tagebuchaufzeichnung hieß das zweite Wort (vergl. die Nachbildung) ursprünglich \textls[120]{plures}; \textsc{Gauss} wollte also etwa schreiben: Warum mehr als \(n\) Wurzeln vorhanden sind u.s.w.}\endgroup, ergibt sich nur aus der \textls[120]{doppelten Periodizität} der lemniskatischen Funktion. \textsc{Gauss} hatte also am 19. März 1797 bereits die doppelte Periode entdeckt. Damit war ihm der Weg geöffnet zur Auffindung der Nullstellen und Unendlichkeitsstellen von sinus und cosinus lemniscaticus und zur Darstellung dieser Funktionen als Quotienten doppelt unendlicher Produkte; siehe die oben S. 153 ff. in den art. [2.]—[6.] abgedruckten Leisteaufzeichnungen. Ferner fand er sich genötigt, seine Funktionen auch für komplexe Werte des Arguments zu betrachten. In der Tat findet sich auf dem letzten Blatte des \textit{Tagebuchs}, mit dem etwas verwischten, aber noch deutlich erkennbaren Datum »1797, Apr. 15.«, also im unmittelbaren Anschluß an die in den Tagebuchnotizen Nr. 60 bis Nr. 63 angezeigten Entdeckungen, die Aufzeichnung:
\par\endgroup
\indent \textls[120]{Quantitates imaginariae:} Quaeritur criterium generale, secundum quod\linebreak functiones plurium variabilium complexae ab incomplexis dignosci possint.

\begingroup\small
\indent Zur Lemniskatenteilung selbst vergl. die oben S. 160 ff. im Abschnitt [III.] zusammengestellten Aufzeichnungen, ferner die Tagebuchnotiz Nr. 62, und die berühmte Stelle im art. 335. des siebenten Abschnitts der \textit{Disquisitiones arithmeticae} (1801), Werke I, S. 412, 413.
\par\noindent\hfill\textsc{Schlesinger.}
\par\endgroup
